\begin{abstract}
Accurate photographic species recognition is vital for biodiversity monitoring, yet fine-grained visual similarities pose steep challenges under limited data. We benchmark vision pipelines on a 1{,}000-species subset of iNaturalist-2021 (40 training images/class), comparing traditional handcrafted baselines (Random Forests with Colour Histograms/HOG, and Linear SVMs with BoVW-RootSIFT) against modern deep CNNs (EfficientNetV2 and ConvNeXt V2). Deep models achieve dramatic gains (ConvNeXt V2 reaching 81.00\% Top-1 accuracy; EfficientNetV2 79.20\%) over handcrafted baselines ($\le 5.14\%$), establishing the necessity of transfer learning. Finally, post-hoc Grad-CAM analyses confirm models focus on authentic anatomical traits rather than background artifacts, while corruption benchmarks reveal notable vulnerability to motion blur despite strong resilience to JPEG compression.
\end{abstract}